\chapter{The $\varepsilon$-regularity theorem}\label{ch:eps-reg}
\module{NoCompromise.Regularity.ApproxHarmonic,
        NoCompromise.Regularity.EpsReg,
        NoCompromise.Regularity.Excess,
        NoCompromise.Regularity.FluxDefect,
        NoCompromise.Regularity.GeometricRecurrence,
        NoCompromise.Regularity.GraphApprox,
        NoCompromise.Regularity.GraphBadBase,
        NoCompromise.Regularity.GraphGoodBase,
        NoCompromise.Regularity.GraphProjectedExcess,
        NoCompromise.Regularity.GraphSingleValued,
        NoCompromise.Regularity.GraphTwoPointCaps,
        NoCompromise.Regularity.HarmonicAffine,
        NoCompromise.Regularity.HarmonicBlowup,
        NoCompromise.Regularity.HarmonicBlowupUniform,
        NoCompromise.Regularity.HeightBound,
        NoCompromise.Regularity.ReversePoincare,
        NoCompromise.Regularity.TiltImprovement,
        NoCompromise.Regularity.ZeroExcess}

\section{Excess and its elementary comparisons}\label{sec:excess}

\begin{definition}[Cylindrical and spherical excess]\label{def:excess}
\uses{not:cylinders,thm:degiorgi-structure}
For $E$ of locally finite perimeter, $x\in\R^3$, $r>0$ and $\nu\in\Sph^2$,
\begin{align}
  \Exc(E,x,r,\nu)
  &:=\frac1{r^2}\int_{\rbd E\cap C_r^\nu(x)}|\nu_E-\nu|^2\,d\Hs^2 ,
  \label{eq:excess-cyl}\\
  \ExcB(E,x,r,\nu)
  &:=\frac1{r^2}\int_{\rbd E\cap B_r(x)}|\nu_E-\nu|^2\,d\Hs^2 .
  \label{eq:excess-ball}
\end{align}
\end{definition}

\begin{lemma}[Ball--cylinder comparison]\label{lem:excess-ball-cyl}
\uses{def:excess}
From $B_r(x)\subset C_r^\nu(x)\subset B_{\sqrt2r}(x)$,
\begin{equation}\label{eq:excess-ball-cyl}
  \ExcB(E,x,r,\nu)\le\Exc(E,x,r,\nu)\le2\,\ExcB(E,x,\sqrt2r,\nu).
\end{equation}
\end{lemma}

\begin{proof}
Use $B_r(x)\subset C_r^\nu(x)\subset B_{\sqrt2r}(x)$.  The second inclusion
changes the normalization from $r^{-2}$ to $(\sqrt2r)^{-2}$ and therefore
produces the factor $2$.
\end{proof}

\begin{lemma}[Change of centre]\label{lem:excess-center}
\uses{def:excess}
If $z\in C_{r/8}^\nu(x)$ then $C_{r/8}^\nu(z)\subset C_r^\nu(x)$, hence
\begin{equation}\label{eq:excess-center}
  \Exc(E,z,r/8,\nu)\le64\,\Exc(E,x,r,\nu).
\end{equation}
\end{lemma}

\begin{proof}
The inclusion $C_{r/8}^\nu(z)\subset C_r^\nu(x)$ bounds the numerator by
that defining $\Exc(E,x,r,\nu)$, while $(r/8)^{-2}=64r^{-2}$.
\end{proof}

\begin{fnote}[These are the only two]\label{fnote:excess-comparisons}
Lemmas~\ref{lem:excess-ball-cyl} and \ref{lem:excess-center} are the only
ball--cylinder and change-of-centre comparisons used in this chapter.  In
particular \eqref{eq:excess-center} is what makes the iteration of
\S\ref{sec:eps-completion} available at \emph{every} nearby boundary point,
not just at the original centre.
\end{fnote}

\section{Zero-excess classification}

\begin{lemma}[Zero-excess classification]\label{lem:zero-excess}
\uses{def:excess,thm:polar}
Let $F$ have finite perimeter in a connected open cylinder $C$ and suppose
\begin{equation}\label{eq:zero-excess-hyp}
  \nu_F=e_3\qquad|D\ind F|\text{-a.e.\ in }C .
\end{equation}
Then, after changing $F$ on a null set, exactly one of the following holds in
$C$: $F\cap C=\varnothing$; $F\cap C=C$; or
$F\cap C=C\cap\{x_3<c\}$ for some $c\in\R$, where the cutting plane is allowed
to lie outside $C$.  In the nonconstant case only the \emph{lower} halfspace
occurs.
\end{lemma}

\begin{proof}\uses{thm:polar,lem:bv-slices,conv:outward}
The polar decomposition \eqref{eq:polar-sign} and
\eqref{eq:zero-excess-hyp} give
\begin{equation}\label{eq:zero-excess-derivs}
  D_i\ind F=0\ (i=1,2),
  \qquad
  D_3\ind F=-|D\ind F| .
\end{equation}
Testing the first two identities against tensor-product test functions and
using Fubini shows $\ind F(x',x_3)=g(x_3)$ a.e.\ in every smaller cylinder.
The last identity says the distributional derivative of the
$\{0,1\}$-valued $g$ is nonpositive, so $g$ is non-increasing; a
non-increasing function taking only the values $0,1$ a.e.\ is identically
$0$, identically $1$, or $\ind{(-\infty,c)}$.  Compatibility on overlapping
smaller cylinders gives the stated form on $C$, including the two degenerate
cases when the cutting plane lies outside the cylinder.
\end{proof}

\begin{fnote}[Both degenerate cases are needed]\label{fnote:zero-excess-degenerate}
The two constant cases are not decoration: they occur in the compactness
argument of Lemma~\ref{lem:height-bound}, where the limit may a priori be a
full or empty cylinder, and they are excluded there by the density estimate at
the origin rather than by hypothesis.  Similarly the specification
\emph{lower} halfspace, rather than ``a halfspace'', is what makes the
orientation in \eqref{eq:caps} determinate.
\end{fnote}

\section{The height bound}

\begin{constants}\label{const:height}
In Lemma~\ref{lem:height-bound} below, $\varepsilon_h$ depends only on $\eta$; the order of choice is
$\eta\rightsquigarrow\varepsilon_h$.
\end{constants}

\begin{lemma}[Height bound]\label{lem:height-bound}
\uses{def:excess,def:omega-minimal,conv:representative}
For every $\eta\in(0,1/8)$ there is $\varepsilon_h(\eta)>0$ such that: if $E$
is $\omega$-minimal, $0\in\partial\Omega$, $0<r\le1$, and
\begin{equation}\label{eq:height-hyp}
  \Exc(E,0,r,e_3)+\omega r\le\varepsilon_h(\eta),
\end{equation}
then
\begin{equation}\label{eq:height-concl}
  \partial\Omega\cap C_{3r/4}\subset\{|x_3|<\eta r\}.
\end{equation}
\end{lemma}

\begin{proof}
\uses{lem:ahlfors,thm:bv-compactness,lem:perimeter-measure-convergence,
      lem:zero-excess,lem:two-sided-density,lem:open-representative}
Scale to $r=1$; the rescaled set is $\omega r$-minimal at scales below
$1/r\ge1$.  If the assertion failed there would be sets $E_j$ with
quasiminimality constants $\omega_j\to0$, admissible comparison scales at
least one, excess tending to zero, and points
\[
  x_j\in\partial\Omega_j\cap C_{3/4},\qquad|(x_j)_3|\ge\eta .
\]
The Ahlfors bounds \eqref{eq:ahlfors} and Theorem~\ref{thm:bv-compactness}
give, after a subsequence, $\ind{E_j}\to\ind F$ in $L^1(C_{7/8})$, and
every ball compactly contained in $C_{7/8}$ has radius less than one, so
\eqref{eq:pmc-admissible-scales} holds there.
Lemma~\ref{lem:perimeter-measure-convergence}\ref{it:pmc-measures} gives
convergence of the perimeter measures on every smaller cylinder whose boundary
carries no limiting perimeter.  Since
$\int_{\rbd{E_j}\cap C_{7/8}}|\nu_{E_j}-e_3|^2\to0$, fix
$0\le\phi\in C_c(C_{7/8})$.  The polar identity gives
\[
 \int\phi|\nu_{E_j}-e_3|^2\,d|D\ind{E_j}|
 =2\int\phi\,d|D\ind{E_j}|+2\int\phi e_3\cdot D\ind{E_j}.
\]
Perimeter-measure convergence handles the first term, while local $L^1$
convergence gives weak convergence of the distributional derivatives and
handles the second.  The left side tends to zero, hence
$D\ind F=-e_3|D\ind F|$ on $C_{7/8}$, which is exactly
\eqref{eq:zero-excess-hyp}.  Thus Lemma~\ref{lem:zero-excess} makes $F$ a
horizontal halfspace.  Because $0\in\partial\Omega_j$ and both phases have a
uniform density at the origin by \eqref{eq:two-sided-density}, the constant
cases are excluded and the limiting plane passes through zero.

After a further subsequence $x_j\to x_\infty$ with
$|(x_\infty)_3|\ge\eta$.  A ball of radius $c\eta$ centred at $x_j$ lies
strictly in one phase of the limiting halfspace, yet
\eqref{eq:two-sided-density} at the boundary point $x_j$ puts a fixed positive
amount of \emph{each} phase in that ball.  This contradicts the local $L^1$
convergence.
\end{proof}

\section{Phase separation and the reverse Poincar\'e inequality}

\begin{lemma}[Exact flux-defect identity for the excess]\label{lem:flux-defect}
\uses{def:slab-cap,def:excess}
Let $E$ be in slab-and-cap configuration $(1,c,\eta_0)$ with
$\eta_0<1$.  For a.e.\ $s\in(1/2,1)$,
\begin{equation}\label{eq:flux}
  \int_{\rbd E\cap(B_s^2\times(-1,1))}\nu_E\cdot e_3\,d\Hs^2=\pi s^2 ,
\end{equation}
and consequently
\begin{equation}\label{eq:flux-defect}
  \frac12\int_{\rbd E\cap(B_s^2\times(-1,1))}|\nu_E-e_3|^2\,d\Hs^2
  =\Per\bigl(E;B_s^2\times(-1,1)\bigr)-\pi s^2 .
\end{equation}
\end{lemma}

\begin{proof}\uses{thm:gauss-green,def:slab-cap,thm:traces,lem:bv-slices,
                    thm:degiorgi-structure}
For a.e. $s$, slicing by $x'\mapsto|x'|$ gives no perimeter mass on the
vertical wall and the interior and exterior traces there agree.  Apply the
trace product formulas of Theorem~\ref{thm:traces} chartwise to the two caps
and the smooth vertical wall of
$D_s:=B_s^2\times(-1,1)$; the two circular edges are $\Hs^2$-null.  This
identifies the boundary measure of $E\cap D_s$.  Gauss--Green
\eqref{eq:gauss-green} for the constant field $e_3$ now has zero lateral
flux, zero upper-cap contribution, and lower-cap contribution $-\pi s^2$
because the trace of $E$ is respectively unrestricted, zero and one on those
three pieces.  Since the outward normal of the lower cap is $-e_3$,
cancellation of the total flux gives \eqref{eq:flux} with the displayed
positive sign.
Expanding $|\nu_E-e_3|^2=2-2\nu_E\cdot e_3$ and using
$\Hs^2(\rbd E\cap\cdot)=\Per(E;\cdot)$
(Theorem~\ref{thm:degiorgi-structure}) gives \eqref{eq:flux-defect}.
\end{proof}

\begin{fnote}[Where orientation enters]\label{fnote:orientation-enters}
\eqref{eq:flux-defect} is \emph{the} point at which orientation enters the
$\varepsilon$-regularity proof.  An unoriented area Jacobian must not be used
to estimate $|\nu_E-e_3|^2$ pointwise: the flux identity \eqref{eq:flux}
depends on the \emph{signed} cap traces \eqref{eq:caps}, and replacing it by an
unoriented bound loses exactly the cancellation that makes the excess a
\emph{defect}.
\end{fnote}

\begin{lemma}[Reverse Poincar\'e]\label{lem:reverse-poincare}
\uses{def:slab-cap,def:excess,def:omega-minimal,thm:deformation}
There is $\eta_0\in(0,1/2)$, absolute, with the following property.  Let $E$ be
$\omega$-minimal, let $0<\rho\le2^{-1/2}$, and suppose that after a rigid
motion $E$ is in slab-and-cap configuration $(\rho,c,\eta_0)$.  Then
\begin{equation}\label{eq:reverse-poincare}
  \Exc(E,0,\rho/2,e_3)
  \le\frac{C}{\rho^4}\int_{\rbd E\cap C_\rho}(x_3-c)^2\,d\Hs^2
   +C\omega\rho .
\end{equation}
If the two phases in \eqref{eq:caps} are reversed, the same holds with
$-e_3$ in place of $e_3$.
\end{lemma}

\begin{proof}
\uses{lem:flux-defect,thm:deformation,def:omega-minimal,def:slab-cap,
      lem:locality}
Put
\[
 \widetilde E:=\{y:\rho y\in E\},\qquad
 \widetilde c:=c/\rho,\qquad \widetilde\omega:=\rho\omega.
\]
Then $\widetilde E$ has slab-and-cap configuration
$(1,\widetilde c,\eta_0)$.  Scaling \eqref{eq:omega-minimal} directly shows
that $\widetilde E$ satisfies the same comparison with coefficient
$\widetilde\omega$ for every competitor supported compactly in a ball of
radius at most $1/\rho$; this larger admissible radius must be retained after
normalising the geometric scale.

Choose
$\sigma\in(1/2,5/8)$ and $\tau\in(7/8,1)$ so that
\eqref{eq:cylindrical-regular-radii} holds for $\widetilde E$ and
\eqref{eq:flux-defect} holds at $s=\sigma$.  The excluded sets of radii are
null, so such a pair exists.  Apply Theorem~\ref{thm:deformation} to obtain
$F$.  Since
\[
 \sqrt{\tau^2+1}<\sqrt2\le1/\rho,
\]
one can choose $R$ with
$\sqrt{\tau^2+1}<R\le1/\rho$.  The support of
$F\triangle\widetilde E$ lies in the closed cylinder
$\overline{B_\tau^2}\times[-1,1]\Subset B_R$.  Hence the scaled
quasiminimality inequality is applicable.  The construction gives matching
traces on the outer wall and caps, while cylindrical regularity removes
perimeter mass from the outer wall; locality therefore cancels the identical
part outside the cylinder and yields
\[
 \Per\bigl(\widetilde E;B_\tau^2\times(-1,1)\bigr)
 \le \Per\bigl(F;B_\tau^2\times(-1,1)\bigr)
     +\widetilde\omega|F\triangle\widetilde E|.
\]
Now $|F\triangle\widetilde E|\le2\pi\tau^2$, and cylindrical regularity at
$\sigma$ decomposes the first perimeter into its core and annular parts.
Combining this with \eqref{eq:deformation-estimate}, and using
$\tau-\sigma>1/4$, gives
\begin{equation}\label{eq:reverse-unit}
  \Per\bigl(\widetilde E;B_\sigma^2\times(-1,1)\bigr)-\pi\sigma^2
  \le C\int_{\rbd{\widetilde E}\cap C_1}
                 (y_3-\widetilde c)^2\,d\Hs^2(y)
       +C\widetilde\omega .
\end{equation}
By \eqref{eq:flux-defect} at $s=\sigma$ and
$C_{1/2}\subset B_\sigma^2\times(-1,1)$,
\[
 \Exc(\widetilde E,0,1/2,e_3)
 \le C\int_{\rbd{\widetilde E}\cap C_1}
              (y_3-\widetilde c)^2\,d\Hs^2(y)
      +C\widetilde\omega .
\]
Finally, cylindrical excess is scale invariant, the displayed height
integral equals
$\rho^{-4}\int_{\rbd E\cap C_\rho}(x_3-c)^2\,d\Hs^2(x)$, and
$\widetilde\omega=\rho\omega$.  This is \eqref{eq:reverse-poincare}.
For the reversed configuration apply the proved case to $E^c$, whose outward
normal is $-\nu_E$ and whose quasiminimality inequality is identical.
\end{proof}

\begin{remark}[The slab hypothesis is supplied, not assumed]
\label{rem:slab-supplied}
Under the hypotheses of Theorem~\ref{thm:eps-regularity},
\eqref{eq:slab} follows from Lemma~\ref{lem:height-bound}, and
\eqref{eq:caps} follows from the fact that the limiting halfspace in
Lemma~\ref{lem:height-bound} has outward normal $e_3$.  When
Lemma~\ref{lem:reverse-poincare} is applied at the \emph{new} affine plane in
Lemma~\ref{lem:tilt-improvement}, verifying \eqref{eq:slab}--\eqref{eq:caps}
requires a separate argument; see \eqref{eq:chebyshev}.  Skipping that
verification is the most common way to produce a false proof of
\eqref{eq:tilt-improvement}.
\end{remark}

\section{Maximal-function graph approximation}\label{sec:graph-approx}

\begin{definition}[Projected excess measure and its maximal density]
\label{def:projected-excess}
\uses{def:excess,not:cylinders,def:maximal}
For $E$ of locally finite perimeter let $\pi(x',x_3):=x'$ and put, for Borel
$A\subset\R^2$,
\begin{equation}\label{eq:projected-excess}
  \mu(A):=\int_{\rbd E\cap\pi^{-1}(A)\cap C_{3/4}}|\nu_E-e_3|^2\,d\Hs^2 ,
\end{equation}
and $M_{1/8}\mu(x'):=\sup_{0<\rho<1/8}\mu(B^2_\rho(x'))/(\pi\rho^2)$.
\end{definition}

\begin{lemma}[Good base set]\label{lem:good-base}
\uses{def:projected-excess,lem:maximal,def:excess}
Let $E$ have locally finite perimeter, let $\mu$ be its projected excess
measure from Definition~\ref{def:projected-excess}, and let $\gamma,c>0$.
Put $c_\gamma:=c\gamma^2$.  Then
\begin{equation}\label{eq:good-base-measure}
  \bigl|\{M_{1/8}\mu>c_\gamma\}\bigr|\le Cc_\gamma^{-1}\mu(B^2_{3/4}),
\end{equation}
so with $G_0:=\{M_{1/8}\mu\le c\gamma^2\}\cap B^2_{1/2}$,
\begin{equation}\label{eq:good-base-complement}
  \bigl|B^2_{1/2}\setminus G_0\bigr|
  \le Cc^{-1}\gamma^{-2}\Exc(E,0,1,e_3).
\end{equation}
\end{lemma}

\begin{proof}\uses{lem:maximal,def:projected-excess}
\eqref{eq:good-base-measure} is Lemma~\ref{lem:maximal}; then
$\mu(B^2_{3/4})\le\Exc(E,0,1,e_3)$ by
Definitions~\ref{def:excess} and \ref{def:projected-excess}.
\end{proof}

\begin{constants}\label{const:graph}
In Lemma~\ref{lem:two-point-lipschitz} below, The order of choice is
$\gamma\rightsquigarrow\tau\rightsquigarrow\eta\rightsquigarrow(c,\varepsilon_g)$:
first the target Lipschitz constant $\gamma$; then a slab width
$\tau<\min\{1/64,\gamma/256\}$; then a height parameter
$\eta:=\min\{\gamma/4,1/4\}$ for Lemma~\ref{lem:height-bound}; then $c$ and
$\varepsilon_g$ small enough that
$Cc\gamma^2\le\varepsilon_h(\eta)$ and
$\varepsilon_g\le\varepsilon_h(\eta)$, and also
$\varepsilon_g|B_1|^{1/3}\le c_I/2$ so that
Constants~\ref{const:rd} gives $r_d=1$ throughout this unit-scale argument.
\end{constants}

\begin{lemma}[Two-point Lipschitz estimate on the good base]
\label{lem:two-point-lipschitz}
\uses{lem:good-base,lem:height-bound,def:excess,lem:ahlfors,
      thm:gauss-green,thm:degiorgi-structure}
Let $\gamma\in(0,1/8)$.  There are $c=c(\gamma)>0$ and
$\varepsilon_g(\gamma)>0$, chosen in the order of Constants~\ref{const:graph},
such that, if $E$ is $\omega$-minimal with
$0\in\partial\Omega$ and
\begin{equation}\label{eq:graph-hyp}
  \Exc(E,0,1,e_3)+\omega\le\varepsilon_g(\gamma),
\end{equation}
then any two reduced-boundary points $p=(p',p_3)$, $q=(q',q_3)$ with
$p,q\in C_{1/2}$ and $p',q'\in G_0$ satisfy
\begin{equation}\label{eq:two-point-lipschitz}
  |p_3-q_3|\le\gamma|p'-q'| .
\end{equation}
Moreover the phases on the two horizontal faces are fixed, in the trace
sense:
\begin{equation}\label{eq:graph-cap-separation}
 B_{1/2}^2\times\{-1/2\}\subset E^{(1)},
 \qquad
 B_{1/2}^2\times\{1/2\}\subset E^{(0)}
 \quad\text{up to }\Hs^2\text{-null sets}.
\end{equation}
\end{lemma}

\begin{proof}
\uses{lem:height-bound,lem:good-base,def:projected-excess,const:graph,
      lem:ahlfors,thm:gauss-green,thm:degiorgi-structure}
First apply Lemma~\ref{lem:height-bound} at $(0,1)$ with a slab width
\begin{equation}\label{eq:slab-width-choice}
  \tau<\min\{1/64,\ \gamma/256\},
\end{equation}
so that $\partial\Omega\cap C_{3/4}\subset\{|x_3|<\tau\}$.  Since
$D\ind E$ vanishes in each of the connected regions
$B_{3/4}^2\times(-3/4,-\tau)$ and
$B_{3/4}^2\times(\tau,3/4)$, the two corresponding phase values are constants
$\ell,u\in\{0,1\}$.  For a.e. $s\in(1/2,5/8)$, Gauss--Green in
$B_s^2\times(-1/2,1/2)$ gives
\[
 \int_{\rbd E\cap(B_s^2\times(-1/2,1/2))}\nu_E\cdot e_3\,d\Hs^2
 =\pi s^2(\ell-u).
\]
On the other hand the left side is the perimeter in that cylinder minus one
half of its excess numerator.  The cylinder contains $B_{1/4}$, so the
Ahlfors lower bound makes its perimeter at least a fixed positive constant,
whereas \eqref{eq:graph-hyp} makes the excess numerator arbitrarily small.
Thus the left side is positive, forcing $(\ell,u)=(1,0)$.  This proves
\eqref{eq:graph-cap-separation}; in particular no unrecorded replacement of
$E$ by its complement is being made.

Set $d:=|p'-q'|$ and $h:=|p_3-q_3|$.  If $d\ge1/32$ then the global slab
bound and \eqref{eq:slab-width-choice} give $h\le2\tau<\gamma d$.  Suppose
$d<1/32$ and, towards a contradiction, $h>\gamma d$.  Put
$\rho:=2\max\{d,h\}$; the global slab bound gives $\rho<1/8$ and
$q\in C_{3\rho/4}(p)$.  Since $p'\in G_0$,
Definition~\ref{def:projected-excess} gives
\begin{equation}\label{eq:excess-at-p}
  \Exc(E,p,\rho,e_3)\le Cc\gamma^2 .
\end{equation}
By Constants~\ref{const:graph}, Lemma~\ref{lem:height-bound} applies at
$(p,\rho)$ with height parameter $\eta=\min\{\gamma/4,1/4\}$ and gives
$h<\eta\rho$.  If $d\ge h$ this says $h<2\eta d\le\tfrac\gamma2d$,
contradicting $h>\gamma d$.  If $h>d$ it says $h<2\eta h\le h/2$, again a
contradiction.  The case $d=0$ is included: it forces $h=0$.
\end{proof}

\begin{lemma}[Single-valuedness from slicing]\label{lem:single-valued}
\uses{lem:two-point-lipschitz,lem:bv-slices}
In the situation of Lemma~\ref{lem:two-point-lipschitz}, for a.e.\
$x'\in B^2_{1/2}$ the restricted vertical slice
$t\mapsto\ind E(x',t)$ on $(-1/2,1/2)$ is a one-dimensional BV function
whose oriented jump number between its lower and upper traces is exactly one.
If $x'\in G_0$, \eqref{eq:two-point-lipschitz} permits at most one
reduced-boundary point in $\{x'\}\times(-1/2,1/2)$.  Hence, outside a planar
null set, $\rbd E\cap C_{1/2}\cap\pi^{-1}(G_0)$ is a single-valued
$\gamma$-Lipschitz graph over a Borel set $G\subset G_0$.
\end{lemma}

\begin{proof}\uses{lem:bv-slices,lem:two-point-lipschitz}
Lemma~\ref{lem:bv-slices} gives the slice BV property.  The slab conclusion
places no boundary on the portions with $|x_3|$ between the thin slab and
$1/2$; equation \eqref{eq:graph-cap-separation} says that their traces are
the lower-one and upper-zero phases.
Thus the one-dimensional fundamental theorem for BV on $(-1/2,1/2)$ makes
the oriented jump number equal to one.  The rest is
\eqref{eq:two-point-lipschitz}.
\end{proof}

\begin{lemma}[Bad-base area bound]\label{lem:bad-base-area}
\uses{lem:single-valued,def:excess}
With $B:=B^2_{1/2}\setminus G$,
\begin{equation}\label{eq:bad-base-area}
\begin{aligned}
  \Hs^2\bigl(\rbd E\cap C_{1/2}\cap\pi^{-1}(B)\bigr)
  &\le|B|+3\int_{\rbd E\cap C_{1/2}\cap\pi^{-1}(B)}
                    (1-\nu_3)\,d\Hs^2\\
  &\le|B|+\tfrac32\Exc(E,0,1,e_3).
\end{aligned}
\end{equation}
\end{lemma}

\begin{proof}\uses{lem:single-valued,def:excess}
Write $1=(1-|\nu_3|)+|\nu_3|$ and integrate over
$S_B:=\rbd E\cap C_{1/2}\cap\pi^{-1}(B)$.  The localized signed slicing
identity from Lemma~\ref{lem:single-valued} says that the signed sum of jump
orientations in $(-1/2,1/2)$ is one, and therefore bounds
$\int_{S_B}|\nu_3|\,d\Hs^2$ by
$|B|+\int_{S_B}(|\nu_3|-\nu_3)\,d\Hs^2$.  Now use the elementary
inequalities
\begin{equation}\label{eq:elementary-nu3}
  1-|t|\le1-t,
  \qquad
  |t|-t\le2(1-t)
  \qquad(-1\le t\le1),
\end{equation}
and $|\nu_E-e_3|^2=2(1-\nu_3)$.
\end{proof}

\begin{theorem}[Graph approximation]\label{thm:graph-approx}
\uses{lem:two-point-lipschitz,lem:single-valued,lem:bad-base-area}
For every $\gamma\in(0,1/8)$ there are $C=C(\gamma)<\infty$ and
$\varepsilon_g(\gamma)>0$ with the following property; the constant $C$ below
depends only on $\gamma$, and every later use fixes $\gamma$ before any other
constant (Lemmas~\ref{lem:blowup-compactness} and
\ref{lem:harmonic-approx-uniform}, Constants~\ref{const:tilt}), so no
uniformity in $\gamma$ is asserted or needed.  Let $E$ be $\omega$-minimal
with $0\in\partial\Omega$ and \eqref{eq:graph-hyp}.  Then there are a Borel set $G\subset B^2_{1/2}$ and a
function
\begin{equation}\label{eq:graph-lipschitz}
  f:B^2_{1/2}\to\R,\qquad\Lip(f)\le\gamma ,
\end{equation}
such that
\begin{equation}\label{eq:graph-identity}
  \rbd E\cap C_{1/2}\cap\pi^{-1}(G)
  =\{(x',f(x')):x'\in G\}
\end{equation}
up to $\Hs^2$-null sets,
\begin{equation}\label{eq:graph-loss}
  \bigl|B^2_{1/2}\setminus G\bigr|
  +\Hs^2\bigl((\rbd E\cap C_{1/2})\setminus\operatorname{graph}f\bigr)
  \le C\gamma^{-2}\Exc(E,0,1,e_3),
\end{equation}
and
\begin{equation}\label{eq:graph-dirichlet}
  \int_{B^2_{1/2}}|\nabla f|^2\le C\,\Exc(E,0,1,e_3).
\end{equation}
\end{theorem}

\begin{proof}
\uses{lem:two-point-lipschitz,lem:single-valued,lem:bad-base-area,
      lem:good-base,cor:graph-area}
Lemma~\ref{lem:single-valued} produces the single-valued
$\gamma$-Lipschitz graph over $G\subset G_0$; extend its height from $G$ to all
of $B^2_{1/2}$ by the McShane formula, which preserves the Lipschitz constant.
If the height bound used in Lemma~\ref{lem:two-point-lipschitz} is
$|x_3|<\tau$, compose this extension with the nearest-point projection onto
$[-\tau,\tau]$.  This projection is $1$-Lipschitz and fixes the values on
$G$, so the resulting extension still has Lipschitz constant at most
$\gamma$, has the same good graph, and also satisfies
\begin{equation}\label{eq:graph-extension-height}
 \|f\|_{L^\infty(B^2_{1/2})}\le\tau.
\end{equation}
This gives \eqref{eq:graph-lipschitz}--\eqref{eq:graph-identity}.

Equation \eqref{eq:graph-loss} combines
\eqref{eq:good-base-complement} with \eqref{eq:bad-base-area}.

For \eqref{eq:graph-dirichlet}: on the good graph,
\begin{equation}\label{eq:normal-gradient-comparison}
  |\nu_E-e_3|^2=2\Bigl(1-\frac1{\sqrt{1+|\nabla f|^2}}\Bigr)\ge c|\nabla f|^2
\end{equation}
because $|\nabla f|\le\gamma\le1/8$.  Here the signed slice identity and the
lower-one/upper-zero cap separation orient every good jump upward, so the
outward normal is $(-\nabla f,1)/\sqrt{1+|\nabla f|^2}$; without this sign
information the first equality would be false.  On the bad planar set the McShane
extension has $|\nabla f|\le\gamma$ a.e., so its Dirichlet energy there is at
most $\gamma^2|B|\le C\Exc(E,0,1,e_3)$ by \eqref{eq:good-base-complement}.
Corollary~\ref{cor:graph-area} on the good graph converts
\eqref{eq:normal-gradient-comparison} into
\eqref{eq:graph-dirichlet}.
\end{proof}

\section{Approximate harmonicity}

\begin{lemma}[Approximate harmonicity]\label{lem:approx-harmonic}
\uses{thm:graph-approx,lem:bounded-H}
Fix $\gamma\in(0,1/8)$.  There is $C_\gamma<\infty$, depending only on
$\gamma$ and chosen before the clamping height in
\eqref{eq:graph-extension-height}, such that the graph approximation $f$ of
Theorem~\ref{thm:graph-approx} may be taken to satisfy in addition, for every
$\zeta\in C^1_c(B^2_{1/2})$,
\begin{equation}\label{eq:approx-harmonic-clean}
  \Bigl|\int_{B^2_{1/2}}\nabla f\cdot\nabla\zeta\Bigr|
  \le C_\gamma\bigl(\Exc(E,0,1,e_3)+\omega\bigr)\|\nabla\zeta\|_\infty .
\end{equation}
At a general scale $r\in(0,1]$, for the graph approximation of the set
rescaled to unit scale,
\begin{equation}\label{eq:approx-harmonic-scaled}
  \Bigl|\int_{B^2_{r/2}}\nabla f\cdot\nabla\zeta\Bigr|
  \le C_\gamma\,r^2\bigl(\Exc(E,0,r,e_3)+\omega r\bigr)\|\nabla\zeta\|_\infty .
\end{equation}
\end{lemma}

\begin{proof}
\uses{lem:bounded-H,thm:graph-approx,lem:height-bound,lem:ahlfors,
      cor:graph-area,const:graph}
Let $\eta\in C^1_c((-3/4,3/4))$ be identically one on the slab containing
$\partial\Omega\cap C_{3/4}$ (available by
Lemma~\ref{lem:height-bound}), and use in \eqref{eq:bounded-H} the vertical
field $X(x',x_3):=\zeta(x')\eta(x_3)e_3$; the derivatives of $\eta$ vanish on
the boundary portion under consideration.  Since
$\spt\zeta\subset B^2_{1/2}$, this field is compactly supported in $B_1$, as
required by Lemma~\ref{lem:bounded-H}.  On the good graph, with
$Q:=\sqrt{1+|\nabla f|^2}$,
\begin{equation}\label{eq:graph-divergence}
  \dvgt X=\frac{\nabla f\cdot\nabla\zeta}{Q^2},
  \qquad
  d\Hs^2=Q\,dx',
\end{equation}
so the graph contribution is $\int_G\nabla f\cdot\nabla\zeta/Q$.  Since
$|\nabla f|\le\gamma$,
\begin{equation}\label{eq:harmonic-err-1}
  \Bigl|\int_G\Bigl(1-\frac1Q\Bigr)\nabla f\cdot\nabla\zeta\Bigr|
  \le C\gamma\|\nabla\zeta\|_\infty\int_G|\nabla f|^2
  \le C\gamma\,\Exc\,\|\nabla\zeta\|_\infty
\end{equation}
by \eqref{eq:graph-dirichlet}.  The boundary omitted by the graph has area
bounded by \eqref{eq:graph-loss}, and on it
$|\dvgt X|\le C\|\nabla\zeta\|_\infty$, so its contribution is at most
\begin{equation}\label{eq:harmonic-err-2}
  C\gamma^{-2}\,\Exc\,\|\nabla\zeta\|_\infty .
\end{equation}
The McShane extension contributes on the bad planar set at most
\begin{equation}\label{eq:harmonic-err-3}
  \int_{B^2_{1/2}\setminus G}|\nabla f||\nabla\zeta|
  \le\gamma\bigl|B^2_{1/2}\setminus G\bigr|\|\nabla\zeta\|_\infty
  \le C\gamma^{-1}\,\Exc\,\|\nabla\zeta\|_\infty .
\end{equation}
Finally the right side of \eqref{eq:bounded-H} is at most
$\omega\int_{\rbd E}|X\cdot\nu_E|\le C\omega\|\zeta\|_\infty$.  Indeed,
cover the portion of $\rbd E\cap B_1$ meeting the support of $X$ by a fixed
dimension-dependent number of balls of radius $1/4$ centred on the boundary,
and apply the Ahlfors upper bound \eqref{eq:ahlfors} in each ball; the
smallness convention in Constants~\ref{const:graph} makes these constants
absolute.  Combining
\eqref{eq:harmonic-err-1}--\eqref{eq:harmonic-err-3} gives the two-term
estimate
\begin{equation}\label{eq:approx-harmonic}
  \Bigl|\int_{B^2_{1/2}}\nabla f\cdot\nabla\zeta\Bigr|
  \le C_\gamma\,\Exc(E,0,1,e_3)\|\nabla\zeta\|_\infty
   +C\omega\|\zeta\|_\infty ,
\end{equation}
with $C_\gamma\le C(1+\gamma^{-2})$.  Since
$\|\zeta\|_\infty\le C\|\nabla\zeta\|_\infty$ for compactly supported
$\zeta$ on $B^2_{1/2}$, this yields \eqref{eq:approx-harmonic-clean}.  Only
the clean form is used later, and only the fact that $C_\gamma$ depends on
$\gamma$ alone matters; the two-term form and the explicit growth in
$\gamma^{-2}$ are not part of the statement.

\emph{Scaling.}  Write $x=r\xi$, $f(x')=r\widetilde f(\xi')$,
$\widetilde\zeta(\xi'):=\zeta(r\xi')$.  The physical left side is $r$ times
the unit-scale one, while
$\|\nabla_\xi\widetilde\zeta\|_\infty=r\|\nabla_x\zeta\|_\infty$; this gives
the factor $r^2$ in \eqref{eq:approx-harmonic-scaled}.
\end{proof}

\section{Harmonic blow-up}\label{sec:harmonic-blowup}

\begin{lemma}[Compactness of the normalised graphs]\label{lem:blowup-compactness}
\uses{thm:graph-approx,lem:approx-harmonic,thm:rellich,
      lem:two-point-lipschitz}
Fix $\gamma\in(0,1/8)$, and use this same graph parameter for every $j$.
Let $E_j$ be $\omega_j$-minimal at unit scale with
$0\in\partial\Omega_j$, satisfying \eqref{eq:graph-hyp}, and put
\begin{equation}\label{eq:blowup-scale}
  a_j:=\bigl(\Exc(E_j,0,1,e_3)+\omega_j\bigr)^{1/2}>0,
  \qquad a_j\longrightarrow0 .
\end{equation}
Let $f_j$ be the graph approximations of Theorem~\ref{thm:graph-approx} and
\begin{equation}\label{eq:blowup-normalized}
  \bar f_j:=\avg_{B^2_{1/2}}f_j,
  \qquad
  g_j:=\frac{f_j-\bar f_j}{a_j}.
\end{equation}
Then, after a subsequence,
\begin{equation}\label{eq:blowup-convergence}
  g_j\rightharpoonup h\ \text{in }H^1(B^2_{1/2}),
  \qquad
  g_j\to h\ \text{in }L^2(B^2_{1/2}),
\end{equation}
and $h$ is harmonic (hence smooth).
\end{lemma}

\begin{proof}
\uses{thm:graph-approx,lem:approx-harmonic,thm:rellich,thm:bv-poincare,
      lem:sobolev-Ck}
\eqref{eq:graph-dirichlet} and Poincar\'e give a uniform $H^1$ bound for
$g_j$; Theorem~\ref{thm:rellich} gives \eqref{eq:blowup-convergence} after a
subsequence.  Dividing \eqref{eq:approx-harmonic-clean} by $a_j$ makes the
right side $O(a_j)$, so $\int\nabla h\cdot\nabla\zeta=0$ for every test
function.  Weakly harmonic implies smooth by
Lemma~\ref{lem:sobolev-Ck}.
\end{proof}

\begin{fnote}[Strong $L^2$, not strong $H^1$]\label{fnote:only-L2}
Only strong $L^2$ convergence is claimed.  Strong $H^1$ would require an
additional energy-convergence argument and is \emph{unnecessary}, because
Lemma~\ref{lem:reverse-poincare} converts height control into tilt control.
Do not attempt to prove the stronger statement.
\end{fnote}

\begin{lemma}[Uniform harmonic approximation]\label{lem:harmonic-approx-uniform}
\uses{lem:blowup-compactness,thm:graph-approx}
Fix $\gamma\in(0,1/8)$.  There is $C_b(\gamma)<\infty$ such that for every
$\tau_0\in(0,1]$ there is
$0<\varepsilon_b(\tau_0,\gamma)\le\varepsilon_g(\gamma)$ with the following
property.  If $E$ is $\omega$-minimal at unit scale with
$0\in\partial\Omega$ and
$a:=\bigl(\Exc(E,0,1,e_3)+\omega\bigr)^{1/2}$ satisfies
$0<a^2\le\varepsilon_b$, then, for every graph approximation $f$ furnished
by Theorem~\ref{thm:graph-approx}, with
$\bar f:=\avg_{B^2_{1/2}}f$, there is a harmonic $h$ on $B^2_{1/4}$ with
\begin{equation}\label{eq:harmonic-approx}
  \int_{B^2_{1/4}}\Bigl|\frac{f-\bar f}{a}-h\Bigr|^2\le\tau_0,
  \qquad
  \int_{B^2_{1/4}}\bigl(|h|^2+|\nabla h|^2\bigr)\le C_b(\gamma).
\end{equation}
The same assertion includes the clamped extension in
\eqref{eq:graph-extension-height}.  The bound $C_b(\gamma)$ is independent
of $\tau_0$ and of the chosen slab width.
\end{lemma}

\begin{proof}
\uses{lem:blowup-compactness,thm:graph-approx,thm:bv-poincare}
The gradient estimate \eqref{eq:graph-dirichlet} and the planar
Poincar\'e inequality give a uniform $H^1$ bound, depending only on the
fixed $\gamma$, for $(f-\bar f)/a$.  Choose $C_b(\gamma)$ larger than the
square of that bound.  If the approximation failed for some $\tau_0$, choose
sets with $a_j\downarrow0$ and a failing graph approximation for each set.
Lemma~\ref{lem:blowup-compactness} gives a harmonic limit $h$ and strong
$L^2$ convergence on $B^2_{1/4}$; weak lower semicontinuity gives its stated
$H^1$ bound.  This contradicts the failure.  Clamping preserves the graph,
Lipschitz and Dirichlet estimates used here, so the bound is uniform over
those extensions as well.
\end{proof}

\section{Improvement of tilt}

\begin{lemma}[Affine approximation of a harmonic function]\label{lem:harmonic-affine}
\uses{lem:harmonic-deriv}
For $h$ harmonic on $B^2_{1/4}$ and $\theta\in(0,1/32)$,
\begin{equation}\label{eq:harmonic-affine}
  \int_{B^2_{4\theta}}\bigl|h-h(0)-\nabla h(0)\cdot x'\bigr|^2
  \le C\theta^6\int_{B^2_{1/4}}|\nabla h|^2 .
\end{equation}
\end{lemma}

\begin{proof}\uses{lem:harmonic-deriv}
Interior derivative estimates \eqref{eq:harmonic-deriv} bound $\|D^2h\|_\infty$
on $B^2_{1/8}$ by $C\|\nabla h\|_{L^2(B^2_{1/4})}$; Taylor's theorem then gives
a pointwise bound $C\theta^2\|D^2h\|_\infty$ on $B^2_{4\theta}$, and squaring
and integrating over a set of area $C\theta^2$ gives $\theta^6$.
\end{proof}

\begin{constants}\label{const:tilt}
Fix the graph parameter once and for all at $\gamma_0:=1/16$.  The order of
choice in Lemma~\ref{lem:tilt-improvement} is
\[
 \gamma_0\rightsquigarrow C\rightsquigarrow\theta
 \rightsquigarrow\tau_t\rightsquigarrow\tau_0
 \rightsquigarrow\varepsilon_t(\theta).
\]
The graph, harmonic and reverse-Poincar\'e bounds at this fixed $\gamma_0$
give $C$ before the scale $\theta$ is chosen.  For each
$\theta\in(0,1/32)$ choose a slab width $\tau_t$, put $\tau_0=\theta^6$,
and then shrink $\varepsilon_t(\theta)$ to meet the explicitly listed
smallness conditions in the proof.  Corollary~\ref{cor:excess-decay} chooses
one scale after this uniform $C$ is known.
\end{constants}

\begin{lemma}[Improvement of tilt]\label{lem:tilt-improvement}
\uses{lem:reverse-poincare,lem:harmonic-approx-uniform,lem:harmonic-affine}
There is an absolute $C\ge1$ such that for every $\theta\in(0,1/32)$ there
is $\varepsilon_t(\theta)>0$ with the following property.  If $E$ is
$\omega$-minimal at unit scale with $0\in\partial\Omega$, put
\begin{equation}\label{eq:tilt-hyp}
  a:=\bigl(\Exc(E,0,1,e_3)+\omega\bigr)^{1/2}
  \quad\text{and assume}\quad a^2\le\varepsilon_t(\theta).
\end{equation}
Then there is a unit vector $\nu'$ with $|\nu'-e_3|\le Ca$ such that
\begin{equation}\label{eq:tilt-improvement}
  \Exc(E,0,\theta,\nu')
  \le C\theta^2\,\Exc(E,0,1,e_3)+C\theta\omega .
\end{equation}
In particular the same $C$ works for all these scales; only the admissible
initial excess $\varepsilon_t(\theta)$ changes with $\theta$.
\end{lemma}

\begin{proof}
\uses{lem:harmonic-approx-uniform,lem:harmonic-affine,lem:reverse-poincare,
      lem:height-bound,thm:graph-approx,lem:two-point-lipschitz,
      lem:harmonic-deriv,def:slab-cap,const:tilt}
Fix $\gamma_0=1/16$.  In this proof every generic constant $C$ depends only
on the estimates for that fixed graph parameter, the harmonic derivative
and affine-approximation estimates, and the absolute constants of
Lemma~\ref{lem:reverse-poincare}.  None depends on $\theta$.
The $H^1$ bound in \eqref{eq:harmonic-approx} and
Lemma~\ref{lem:harmonic-deriv} give a fixed $C_h\ge1$ such that every
harmonic approximant used below satisfies
$|h(0)|+|\nabla h(0)|\le C_h$.

Now let $\theta\in(0,1/32)$.  Choose
\begin{equation}\label{eq:tilt-slab-choice}
 0<\tau_t<\min\{\theta^3,\eta_0\theta/32,1/64,\gamma_0/256\},
 \qquad \tau_0:=\theta^6,
\end{equation}
where $\eta_0$ is the fixed slab constant of
Lemma~\ref{lem:reverse-poincare}.  Choose the positive
$\varepsilon_t(\theta)$ so that
\[
 \varepsilon_t(\theta)\le\min\left\{
 \varepsilon_g(\gamma_0),\ \varepsilon_h(\tau_t),\
 \varepsilon_b(\theta^6,\gamma_0),\ \theta^6,\
 \left(\frac{\eta_0\theta}{64C_h}\right)^2,\
 \left(\frac1{16C_h}\right)^2\right\}.
\]
These choices affect the hypotheses, not the constants in the estimates.
If $a=0$, take $\nu'=e_3$; zero excess at unit scale gives zero excess in
the smaller cylinder.  Assume $a>0$.

The height bound gives $\partial\Omega\cap C_{3/4}\subset\{|x_3|<\tau_t\}$.
Choose the graph approximation with its extension clamped to
$[-\tau_t,\tau_t]$, so $|\bar f|\le\tau_t$.  The phase-separation argument
of Lemma~\ref{lem:two-point-lipschitz} gives the lower-one and upper-zero
regions outside this slab in $C_{3/4}$: each of the two boundary-free
regions is connected, so its phase is constant, and the cap traces in that
lemma identify the two constants.  Shrinking the slab is permitted by
\eqref{eq:tilt-slab-choice}.  Apply
Lemma~\ref{lem:harmonic-approx-uniform} to this graph with tolerance
$\theta^6$, and define
\begin{equation}\label{eq:new-normal}
  \nu':=\frac{(-a\nabla h(0),1)}{\sqrt{1+a^2|\nabla h(0)|^2}}.
\end{equation}
Then $|\nu'-e_3|\le2C_ha\le1/8$.  In particular the rotated cylinder
$C_{2\theta}^{\nu'}$ lies inside $C_{3\theta}\subset C_{3/4}$ and its
horizontal projection lies inside $B^2_{3\theta}\subset B^2_{4\theta}$.

\emph{Step 1: height above the new affine plane.}  Let the plane be
\begin{equation}\label{eq:new-plane}
  x_3=\bar f+a\bigl(h(0)+\nabla h(0)\cdot x'\bigr).
\end{equation}
On the portion of the boundary contained in $\operatorname{graph}f$,
its squared vertical discrepancy is at most
\[
 2a^2\left(\int_{B^2_{1/4}}\left|\frac{f-\bar f}{a}-h\right|^2
 +\int_{B^2_{4\theta}}|h-h(0)-\nabla h(0)\cdot x'|^2\right)
 \le Ca^2\theta^6
\]
after integration over the base.  Here
\eqref{eq:harmonic-approx} and \eqref{eq:harmonic-affine} apply, and the
area element of the graph is bounded by $\sqrt{1+\gamma_0^2}$.
Distance to the affine plane is no larger than this vertical discrepancy.

For a boundary point in $C_{2\theta}^{\nu'}$ outside $\operatorname{graph}f$,
$|x_3|\le\tau_t$, $|\bar f|\le\tau_t$ and the bound on $h(0),\nabla h(0)$
make its squared distance at most $C(\tau_t^2+a^2)$.  By
\eqref{eq:graph-loss}, the area of such points is bounded by
$C\gamma_0^{-2}\Exc(E,0,1,e_3)$.  Since
$\tau_t^2\le\theta^6$, $a^2\le\theta^6$ and
$\Exc(E,0,1,e_3)\le a^2$, their contribution is at most $Ca^2\theta^6$.
Thus the full boundary height-square in the rotated cylinder is at most
$Ca^2\theta^6$, with $C$ independent of $\theta$.

\emph{Step 2: slab and caps relative to the new plane.}
Write $b:=\bar f+ah(0)$ and
$c:=b/\sqrt{1+a^2|\nabla h(0)|^2}$, the signed distance of the plane from
zero.  The choices above give
\begin{equation}\label{eq:chebyshev}
 |b|+3\theta a|\nabla h(0)|
 \le\tau_t+(1+3\theta)C_ha<\eta_0\theta/8.
\end{equation}
Consequently every boundary point in $C_{2\theta}^{\nu'}$ has distance
less than $\tau_t+\eta_0\theta/8<2\eta_0\theta$ from the new plane.
Also $|c|\le|b|<\eta_0\theta/8$, so
$|c|+2\eta_0\theta<2\theta$ because $\eta_0<1/2$.
This verifies both the slab and clearance conditions of
Definition~\ref{def:slab-cap} at radius $2\theta$.

A point on either cap has the form $x=\pm2\theta\nu'+y$, with
$y\cdot\nu'=0$ and $|y|<2\theta$.  Since $|\nu'-e_3|\le1/8$, its third
coordinate is less than $-\theta$ on the lower cap and greater than
$\theta$ on the upper cap.  Both caps lie in $C_{3/4}$ and outside the
original slab, so their traces are respectively one and zero by the phase
separation already established.  The complete slab-and-cap configuration
therefore holds after rotating $\nu'$ to $e_3$.

\emph{Step 3: apply reverse Poincar\'e.}
Since $2\theta<2^{-1/2}$, Lemma~\ref{lem:reverse-poincare} applies and gives
\[
  \Exc(E,0,\theta,\nu')
  \le C\theta^{-4}\bigl(a^2\theta^6\bigr)+C\omega\theta.
\]
Substitute $a^2=\Exc(E,0,1,e_3)+\omega$ and use $\theta<1$ to absorb
$C\theta^2\omega$ into $C\theta\omega$.  Every constant in these three
steps was independent of $\theta$; choose one common $C\ge1$ larger than
those constants and $2C_h$.  This establishes the quantified statement.
\end{proof}

\begin{corollary}[Excess decay]\label{cor:excess-decay}
\uses{lem:tilt-improvement}
Let $C\ge1$ be the uniform constant in Lemma~\ref{lem:tilt-improvement}, and
choose once and for all
\begin{equation}\label{eq:theta-choice}
 \theta:=\min\{1/64,1/(4C)\}\in(0,1/32),
 \qquad C\theta^2\le\theta/4\le\theta/2.
\end{equation}
Then there is $\varepsilon_0>0$, absolute, such that: if $E$ is
$\omega$-minimal, $0\in\partial\Omega$, $0<r\le1$, and
\begin{equation}\label{eq:eps-reg-hyp}
  \Exc(E,0,r,e_3)+\omega r\le\varepsilon_0,
\end{equation}
then there is a unit vector $\nu'$ with
\begin{equation}\label{eq:excess-decay}
  \Exc(E,0,\theta r,\nu')
  \le\frac\theta2\,\Exc(E,0,r,e_3)+C\theta\omega r.
\end{equation}
\end{corollary}

\begin{proof}\uses{lem:tilt-improvement}
Obtain $C$ first from Lemma~\ref{lem:tilt-improvement}, then choose
$\theta$ by \eqref{eq:theta-choice}, then set
$\varepsilon_0:=\varepsilon_t(\theta)$.  Under $x=r\xi$ the excess is
invariant and the quasiminimality coefficient becomes $\omega r$, with
admissible comparison radii at least one since $r\le1$.  Apply the lemma
at this scale and use $C\theta^2\le\theta/2$.
\end{proof}

\section{Iteration and completion}\label{sec:eps-completion}

\begin{lemma}[Geometric recurrence]\label{lem:geometric-recurrence}
Let $E_0\ge0$, $\theta\in(0,1)$, $\omega,r\ge0$, and let $E_j\ge0$ satisfy
\begin{equation}\label{eq:recurrence}
  E_{j+1}\le\frac\theta2E_j+C\theta\omega r_j,
  \qquad r_j:=\theta^jr .
\end{equation}
Then $E_j\le C'\theta^j(E_0+\omega r)$ for every $j$, with $C'$ depending only
on $C$.
\end{lemma}

\begin{proof}
Divide by $\theta^{j+1}$: the recurrence for $\widetilde E_j:=E_j/\theta^j$
becomes $\widetilde E_{j+1}\le\tfrac12\widetilde E_j+C\omega r$, whose
solution is bounded by $\widetilde E_0+2C\omega r$.
\end{proof}

\begin{lemma}[The optimal normals are Cauchy]\label{lem:normals-cauchy}
\uses{cor:excess-decay,lem:ahlfors}
In the iteration of Corollary~\ref{cor:excess-decay} at scales
$r_j=\theta^jr$, with $E_j$ the excess relative to the normal $\nu_j$ selected
at step $j$,
\begin{equation}\label{eq:normal-increment}
  |\nu_{j+1}-\nu_j|^2\le C\bigl(E_j+\omega r_j\bigr) ,
\end{equation}
and consequently $\nu_j\to\nu_\infty$ with
\begin{equation}\label{eq:normal-rate}
  |\nu_j-\nu_\infty|\le C\theta^{j/2}\bigl(E_0+\omega r\bigr)^{1/2}.
\end{equation}
\end{lemma}

\begin{proof}\uses{lem:ahlfors,lem:geometric-recurrence,cor:excess-decay}
On a cylinder of radius comparable to $r_{j+1}$, expand
$|\nu_{j+1}-\nu_j|^2$ through $\nu_E$, bound the two error integrals by the
excesses $E_j$ and $E_{j+1}$, and use the Ahlfors \emph{lower} bound
\eqref{eq:ahlfors} to divide by a positive multiple of $r_{j+1}^2$.
Because $\theta$ is fixed, this gives
$|\nu_{j+1}-\nu_j|^2\le C(E_j+E_{j+1})$; the recurrence
\eqref{eq:recurrence} gives \eqref{eq:normal-increment}.  Then
Lemma~\ref{lem:geometric-recurrence} gives
$E_j+\omega r_j\le C\theta^j(E_0+\omega r)$, and summing
\eqref{eq:normal-increment} gives \eqref{eq:normal-rate}.
\end{proof}

\begin{lemma}[Uniform iteration at nearby centres]\label{lem:uniform-centers}
\uses{lem:excess-center,cor:excess-decay,lem:normals-cauchy}
There is $\varepsilon_0'>0$, absolute, such that if $E$ is $\omega$-minimal,
$0\in\partial\Omega$, $0<r\le1$, and
$\Exc(E,0,r,e_3)+\omega r\le\varepsilon_0'$, then
\eqref{eq:eps-reg-hyp} holds at every $z\in\partial\Omega\cap C_{r/8}$ at
scale $r/8$, and hence the iteration of
Corollary~\ref{cor:excess-decay} and Lemma~\ref{lem:normals-cauchy} runs at
every such centre.
\end{lemma}

\begin{proof}\uses{lem:excess-center}
$C_{r/8}(z)\subset C_{r/2}$, so \eqref{eq:excess-center} gives
$\Exc(E,z,r/8,e_3)\le64\,\Exc(E,0,r,e_3)$; take
$\varepsilon_0':=\varepsilon_0/64$ (also shrinking to absorb the $\omega r/8$
term).
\end{proof}

\begin{theorem}[$\varepsilon$-regularity]\label{thm:eps-regularity}
\uses{cor:excess-decay,lem:normals-cauchy,lem:uniform-centers,conv:representative}
There are absolute constants $\varepsilon_0>0$, $\theta\in(0,1/16)$ and
$C<\infty$ such that the following holds.  Let $E$ be $\omega$-minimal with
$0\in\partial\Omega$, let $0<r\le1$, and assume
\begin{equation}\label{eq:eps-reg-hyp-final}
  \Exc(E,0,r,e_3)+\omega r\le\varepsilon_0 .
\end{equation}
Then $\partial\Omega\cap C_{r/4}$ is, after rotation, the graph of a
$C^{1,1/2}$ function, and
\begin{equation}\label{eq:eps-reg-holder}
  |\nu_\Omega(x)-\nu_\Omega(y)|
  \le Cr^{-1/2}|x-y|^{1/2}
    \bigl(\Exc(E,0,r,e_3)+\omega r\bigr)^{1/2}
\end{equation}
for $x,y\in\partial\Omega\cap C_{r/8}$.
\end{theorem}

\begin{proof}
\uses{cor:excess-decay,lem:normals-cauchy,lem:uniform-centers,
      lem:geometric-recurrence,lem:bv-slices,lem:height-bound,
      lem:two-point-lipschitz,lem:open-representative,
      thm:degiorgi-structure}
\emph{Step 1: iterate.}  Apply \eqref{eq:excess-decay} recursively at the
scales $r_j=\theta^jr$.  Lemma~\ref{lem:geometric-recurrence} gives
$E_j\le C\theta^j(E_0+\omega r)$ and Lemma~\ref{lem:normals-cauchy} gives the
limiting normal $\nu_\infty$ with the rate \eqref{eq:normal-rate}.

\emph{Step 2: run the iteration at every nearby centre.}
Lemma~\ref{lem:uniform-centers}.

\emph{Step 3: one graph.}  The limiting normal at every reduced-boundary point
in a sufficiently small common cylinder makes an angle less than $\pi/4$ with
$e_3$.  Indeed, at a reduced-boundary point, differentiation of the excess
measure and \eqref{eq:normal-rate} identify the limiting normal with the
measure-theoretic normal $\nu_E$ there.  Hence
$D_3\ind E=-\nu_E\cdot e_3\,|D\ind E|$ has one sign in the common cylinder.
By BV slicing (Lemma~\ref{lem:bv-slices}), for a.e.\ vertical line the
one-dimensional slice of $\ind E$ is therefore monotone and has at most one
jump.  After rescaling, the phase-separation conclusion
\eqref{eq:graph-cap-separation} of Lemma~\ref{lem:two-point-lipschitz} gives
uniform lower-one and upper-zero regions, and hence exactly one jump for a.e.\
base point.  Let $f(x')$ be that jump height on this full-measure set.

The preceding a.e.\ graph has a unique continuous extension.  To see this
without making an exceptional-slice assumption, take two sequences of
reduced-boundary points converging to topological-boundary points.  Apply the
height bound at the first point at the smallest dyadic iteration scale which
contains the second.  Because all limiting normals make angle less than
$\pi/4$ with $e_3$, the resulting cone estimate says that two sufficiently
near boundary points with the same horizontal projection must coincide, and
that the vertical separation tends to zero with the horizontal separation.
The reduced boundary is dense in $\partial\Omega$ by
Lemma~\ref{lem:open-representative}, so this estimate passes to the limiting
topological-boundary points.  The lower-one and upper-zero regions show that
every vertical segment over the smaller base disk meets the boundary.
Consequently the extended $f$ is continuous and its graph is the entire
topological boundary in the smaller cylinder, not merely an a.e.\ subset.
Local graph extensions agree on overlaps because they represent the same
density-one open set.

\emph{Step 4: H\"older estimate.}  Let $d:=|x-y|$ and choose the largest
iteration scale $s$ with $s\asymp d$.  Take the fixed comparison factor in
$s\asymp d$ large enough that one ball of radius $cs$ centred at a boundary
point lies in both scale-$s$ cylinders.  Expanding the squared difference of
the two scale-$s$ normals through $\nu_E$ on that ball, using the two excess
bounds and dividing by its Ahlfors lower bound, gives
\[
 |\nu_{x,s}-\nu_{y,s}|
 \le C(d/r)^{1/2}
       \bigl(\Exc(E,0,r,e_3)+\omega r\bigr)^{1/2}.
\]
The two tails from these scale-$s$ normals to their limiting normals have the
same bound by \eqref{eq:normal-rate}.  This proves
\eqref{eq:eps-reg-holder} first on the reduced boundary and then everywhere
by continuity.  The continuous graph has this field as its normal; solving
the normal--gradient relation in the cone $\nu\cdot e_3>1/2$ gives
$f\in C^{1,1/2}$.
\end{proof}

